% Transcription — fonds Alexandre Grothendieck, shelfmark 156-6, batch 2 (pages 21-40).
% Established from the facsimile archives/batches/156-6/batch-02.pdf (a copy of
% 156-6.pdf fetched through the project's relay; no cover sheet in the batch,
% so PDF page k = folder page 20+k), per the skill transcribe-grothendieck. The
% folder is ninety-three pages in five batches; this is the second. Batches
% 1, 3-5 and the first draft of the chapter (folder 156-4, GF IV) were being
% transcribed in parallel and could not be read; the notation is fixed here
% from chapter V (folder 156-5) and from these pages alone.
%
% Pass: Opus 5.5 (claude-opus-5-5), 2026-09-26 — first pass, unchecked against the pages by a human.
%
% All twenty leaves are mathematics in his hand and all are transcribed;
% nothing is skipped and nothing is summarised. His own page numbers 21-40
% stand at the head of each leaf and coincide with the archivists'. One date
% in his hand: « 20 juin », underlined, opening page 34 (noted there). Page 38
% is cancelled by one long diagonal stroke and the first paragraph of page 39
% by several; both are transcribed and the cancellation noted once. The
% bottom paragraph of page 23 is likewise cancelled by diagonals.
%
% Numbered runs, recorded in notes: his formula numbers (1.39)-(1.44) (the
% labels on pages 25 and 26 read like « 1.48 »; his own references on pages
% 26-27 to « (1.42) » settle the first; the second is flagged); the axioms of
% an « atelier » At 4 - At 7 (pages 21, 27; At 1 - At 3 lie in batch 1), the
% conditions C 4 - C 6 and C1 - C7 cited from batch 1, and the starred list
% At*1 - At*5 translating them for (𝓜, 𝔉_𝓜) (pages 28-33). Cross-references
% to the first draft: margin of page 30 « cf GF IV page 54 (cor. 2) et p. 55
% (variante), id. pour C3, p. 58 »; page 38 « (cf GF IV, p. 53) ». These are
% to his own pagination of chapter IV (folder 156-4) and are transcribed as
% written.
%
% What the batch is about. A « atelier » is (𝔉, ≤, ≪) — figures, the
% sub-figure order and the refinement order — with 𝓜 ⊂ 𝔉 its multiplicities
% (irreducibles) and F̃ = {X ∈ 𝓜 | X ≤ F}. Pages 21-22: axiom At 4 (G ≪ F iff
% every X ∈ G̃ is ≪ some Y ∈ F̃), so for figures of finite type the atelier is
% recovered from (𝓜, ≤, ≪, compatibility); At 5 (F^X = {Y ∈ F̃ | X ≪ Y} has a
% least element), i.e. Multomb(F) = {Ỹ} is a preadmissible figure in 𝓜; At 6
% (for compatible X, Y, X ≪ Y ⇒ X ≤ Y), with two corollaries; the injection
% (1.39) X ↦ Omb(X) = 𝓜_{≪X}. Pages 23-24: F ≤ G iff Multomb(F) ≤ Multomb(G)
% iff inclusion; compatibility, sup and inf are carried over (with an NB that
% the converse fails for two points), disjointness too, (1.40)
% Multomb(∅) = ∅; then F ≪ G iff Multomb(F) ≪ Multomb(G), proved on pages
% 25-27 via Omb(F) = |Multomb F|, the retraction φ_F : Omb(F) → F̃ and the
% square (1.42); along the way the « raffinement intérieur » X ≪̊ Y (1.43)
% and At 7 (≪̊ transitive). Page 28: Scholie — an atelier is the same as a set
% 𝓜 with a family 𝔉_𝓜 ⊂ Fig(𝓜) ⊂ 𝔓𝔓(𝓜) (1.44); pages 29-33 write the
% conditions At*1 - At*5 on (𝓜, 𝔉_𝓜) (𝓜* the set of all strata, the
% multistrate Φ_A, stability by bounded unions, separation of points), with a
% Théorème-Scholie on page 32 stating the equivalence of categories. Page 34
% (20 June): the map 𝔉 → Fig(ℒ), F ↦ multomb(F), increasing for ≤ and ≪ but
% not injective; Examples 1 A) (triples (X, Σ, Φ) in a space E, piecewise
% linear) and B) (affine simplices (A, σ)), Example 2 (a « pseudodisque »).
% Pages 36-37: definition of an « atelier ensembliste » (a)-c)), morphisms of
% ateliers, « compatible aux strates », sub-ateliers; page 38 (cancelled) and
% pages 39-40: closed parts of an ordered set with bounded sups, « ample » and
% « spécieuse » parts, the irreducibles of a closed part, At 2 inherited, and
% the example 𝔉 = 𝔓(I), where a closed part is a simplicial complex on I and a
% « spécieuse » one is 𝔓(I') — « c'est pas du tout pareil ! ».
%
% Reading hazards fixed once here. 𝔉 (\mathfrak{F}) is his script F for the
% set of figures, 𝓜 (\mathcal{M}) his script M, ℒ (\mathcal{L}) the set of
% lieux, 𝔓 the power set, as in folder 156-5. His compatibility sign, two
% crossed chevrons with a bar, is set \between. The underlined triangle
% (« Y ◁ F », pages 21, 40) is set \trianglelefteq. ≪ with a small circle above
% is set \overset{\circ}{\ll}. « Omb », « Multomb » (multiombre, page 28),
% « Mulstrat » (page 35) and « Fig » are his abbreviations, set in roman. The
% label « At » is written over another letter on pages 21 and 27 (it first
% looks like « AE »); page 38 names the list « At1 - At7 ». Several stars and
% scribbles above 𝓜* on pages 29-30 are struck marks and are not rendered.
% His underlining is set in \emph. The hand is fast drafting; the slanted left
% margins of pages 21, 22, 27, 30, 33, 36, 37 and 40 were turned and read only
% in part. The equation of page 21's slanted note is offered as \uncertain.
%
% Readings flagged and not settled: « nécessaire » (p. 21, first line); the
% sentence on préadmissibilité and C 5 (p. 21); « paramétrisation » and
% « fidèle » (p. 22); the opening of the Corollaire of p. 23; « réduits »
% (p. 24); the sentence « Il semble qu'il faille le prendre en axiome »
% (p. 26); most of page 30's middle paragraph; page 32's closing bracket;
% the NB of page 37 and most of its margins; Example 2 (p. 35). Variances on
% the page and not repaired: « Multomb(X) = {{x}} » where F, G = {x}, {y} are
% meant (p. 24); Corollaire 2 on p. 24 follows an unnumbered Corollaire on
% p. 23 which it calls « corollaire 1 ».
\documentclass[11pt,a4paper]{article}
\input{../preamble/grothendieck}

\folder{156-6}
\batch{2}
\pages{21}{40}
\dating{1986}
\watermark{Édition de démonstration}
\foldertitle{[Chapitre] VI. Analysis situs (deuxième mouture) : notes manuscrites (18-20/06/1986)}

\begin{document}

\page{21}
\note{pagination de sa main 21 à 40 en tête des feuillets, qui coïncide avec celle des archivistes}

question s'il n'existe une semblable \uncertain{nécessaire}. Aussi, finalement, j'ai envie de poser

\textbf{At 4}\note{le sigle « At » est écrit par-dessus une autre lettre (on lit d'abord « AE ») ; les axiomes At 1 à At 3 sont dans le lot 1. L'axiome est marqué d'un double trait vertical dans la marge} Si $F, G \in \mathfrak{F}$, alors $G \ll F$ ssi pour tout $X \in \widetilde{G}$, $\exists\, Y \in \widetilde{F}$ avec $X \ll Y$.

Ainsi, la relation $\ll$ sur $\mathfrak{F}$ se déduit de la relation $\ll$ sur $\mathcal{M}$.\note{une croix de renvoi sous « la » avant « relation » ; son objet n'est pas lu} Donc, si on \ill{} aux \emph{figures} de type fini, $(\mathfrak{F}, \leq, \ll)$ (i.e. \ill{} l'atelier) se reconstitue entièrement à partir de $(\mathcal{M}, \leq, \ll, \between)$.

\marginal{Cor $\forall F \in \mathfrak{F}$, \ill{} \uncertain{$\mathrm{Omb}(F) = |\mathrm{Multomb}(F)|$} \ill{}}\note{note oblique dans la marge gauche, à la hauteur de ce paragraphe, lue en partie ; comparer la formule du début de la page 25}

\textbf{At 5} Soient $X \in \mathcal{M}$, $F \in \mathfrak{F}$, tels que $X \ll F$. Alors \struck{dans} l'ensemble
\[
  F^{X} = \lbrace Y \in \widetilde{F} \mid \underbrace{X \ll Y}_{X \in \mathrm{Omb}(Y)} \rbrace
\]
(lequel est $\neq \emptyset$ par C 4) admet un plus petit élément (pour $\leq$).

En d'autres termes, l'ensemble de parties
\[
  \mathrm{Multomb}(F) = \lbrace \widetilde{Y} \mid Y \in \widetilde{F}, \ Y \trianglelefteq F \rbrace
\]
est une \emph{figure} dans $\mathcal{M}$, \struck{\ill{}} au sens usuel, \struck{que les $\widetilde{Y}^{\circ} = \widetilde{Y} \smallsetminus \ill{}$} qui est \emph{préadmissible}. (C'est \ill{} en tous cas impliqué par \uncertain{C 5}, et lui est équivalent si on admet que la relation (pour $Y, Y' \in \widetilde{F}$) $Y \leq Y'$, qui implique $Y \ll Y'$ donc $\mathrm{Omb}(Y) \subset \mathrm{Omb}(Y')$, lui est en fait équivalente. Cela équivaut à la condition

\textbf{At 6} Si $X, Y \in \mathcal{M}$ et $X \between Y$, alors $X \ll Y \Rightarrow X \leq Y$ (donc $X \ll Y \Leftrightarrow X \leq Y$).

En d'autres termes,

\page{22}
\textbf{Corollaire 1} Soit $F \in \mathfrak{F}$. Alors dans $\widetilde{F} \subset \mathcal{M}$, les relations d'ordre $\leq$ et $\ll$ coïncident.

\textbf{Corollaire 2} Soient $F$ et $G$ deux figures compatibles ($F \between G$). Alors $F \ll G$ ssi $F \leq G$.

Il suffit de prouver $\Longrightarrow$, et par C 4, on est ramené à \add{le} \struck{prouver} voir pour $X$ ($\in \widetilde{F}$) et $Y$ ($\in \widetilde{G}$) dans $\mathcal{M}$, ce qui n'est autre que C 6.

Ainsi, pour toute $F \in \mathfrak{F}$, la \uncertain{paramétrisation} de $\mathrm{Multomb}(F)$ par $\widetilde{F}$ est \uncertain{fidèle} \ill{} et respecte les relations d'ordre \ill{} (\ill{} $\leq$).

\add{Ceci dit}, pour voir que $\mathrm{Multomb}(F)$ est bien une \emph{figure}, il reste à voir que pour $Y \in \widetilde{F}$,
\[
  \widetilde{Y}^{\circ} \overset{\mathrm{def}}{=} \widetilde{Y} \smallsetminus \bigcup_{\substack{Y' \in \widetilde{F} \\ Y' \leq Y}} \widetilde{Y}'
\]
\note{sous la réunion, entre parenthèses : « ou mieux, $Y' \ll Y$, $Y' \neq Y$ »}
est non vide. Mais il n'y a qu'à \struck{prouver} \add{\uncertain{noter}} que $Y \in \widetilde{Y}^{\circ}$ !

\marginal{NB \ill{} C 6, \ill{} $Y' \in \ldots$ \ill{} il \ill{} $Y' \neq Y$, $Y' \leq Y$ \ill{}}\note{note oblique dans la marge gauche, en face de ce passage, lue par fragments seulement}

Par la relation d'ordre $\ll$ sur $\mathcal{M}$, on trouve une application
\[
  \mathcal{M} \hookrightarrow \mathfrak{P}(\mathcal{M}), \qquad X \longmapsto \mathrm{Omb}(X) = \mathcal{M}_{\ll X} = \lbrace Y \in \mathcal{M} \mid Y \ll X \rbrace \tag{1.39}
\]
\note{le numéro (1.39) est écrit par-dessus un autre}
qui est \emph{injective} ($\ll$ relation d'\emph{ordre}, pas seulement préordre) \add{et la relation d'ordre de $\mathcal{M}$ \ill{} induite}. Les figures \struck{\ill{}} $\mathrm{Multomb}(F)$ sont les images, par cette injection, des $\widetilde{F}$ ($F \in \mathfrak{F}$). \struck{\ill{}}

\page{23}
\textbf{Proposition} Soient $F, G \in \mathfrak{F}$. Conditions équivalentes
\begin{itemize}
  \item[(i)] $F \leq G$
  \item[(i bis)] $\widetilde{F} \subset \widetilde{G}$
  \item[(ii)] $\mathrm{Multomb}(F) \leq \mathrm{Multomb}(G)$ i.e. $\mathrm{Multomb}(F)$ est une sous-figure de $\mathrm{Multomb}(G)$
  \item[(ii bis)] $\mathrm{Multomb}(F) \subset \mathrm{Multomb}(G)$,
\end{itemize}

\emph{Dém.} $F \leq G$ \struck{trivial, et} $\Leftrightarrow \widetilde{F} \subset \widetilde{G}$ \struck{trivial} déjà connu \struck{\ill{}} \add{\ill{}} ($\Rightarrow$ trivial, $\Leftarrow$ car $F = \operatorname{Sup}_{\leq} X$ ($X \in \widetilde{F}$), $G = \operatorname{Sup}_{\leq} Y$ ($Y \in \widetilde{G}$)). L'équivalence (i bis) $\Leftrightarrow$ (ii bis) triviale par l'injectivité de (1.39). Or (ii bis) est impliqué par (ii) (tautologique) ; \ill{} l'inverse, \ill{} que (ii bis) implique que $\mathrm{Multomb}(F)$ est fermé dans $\mathrm{Multomb}(G)$ pour l'inclusion. Mais $\widetilde{F} \subset \widetilde{G}$ \struck{\ill{}} \add{donc} $F \leq G$, et il en \ill{} (\ill{} et trivial) que $\widetilde{F}$ est fermé dans $\widetilde{G}$.

\textbf{Corollaire}! \struck{Conditions équivalentes pour} \add{\uncertain{Soient}} $F, G \in \mathfrak{F}$\note{il écrit « $F, G \in \mathcal{M}$ »} tels que $F \between G$, alors \struck{(i) $F \between G$}
\[
  (*) \qquad \mathrm{Multomb}(F) \between \mathrm{Multomb}(G),
\]
\struck{Quand elles sont vérifiées}, on a
\[
  (**) \qquad \begin{cases} \mathrm{Multomb}(F \vee G) = \mathrm{Multomb}(F) \cup \mathrm{Multomb}(G) \\ \mathrm{Multomb}(F \wedge G) = \mathrm{Multomb}(F) \cap \mathrm{Multomb}(G). \end{cases}
\]

La condition (ii) signifie que $\widetilde{F}$ et $\widetilde{G}$ sont fermés dans $\widetilde{F} \cup \widetilde{G} \subset \mathcal{M}$, pour la relation $\ll$ (induite par $\subset$ dans $\mathfrak{P}(\mathcal{M})$).\note{ce dernier paragraphe est barré de plusieurs traits obliques}

\page{24}
En effet, $\widetilde{F \vee G} = \widetilde{F} \cup \widetilde{G}$, et $\widetilde{F}$ et $\widetilde{G}$ y sont fermés, \struck{dans $\mathcal{M}$} ce qui prouve que
\[
  \mathrm{Multomb}(F \vee G) = \mathrm{Multomb}(F) \cup \mathrm{Multomb}(G),
\]
et que $\mathrm{Multomb}(F)$ et $\mathrm{Multomb}(G)$ y sont fermés, donc (*) et la première relation (**). La deuxième relation provient de ce que $\widetilde{F \wedge G} = \widetilde{F} \cap \widetilde{G}$.

\textbf{NB} Il n'est pas vrai en général que $\mathrm{Multomb}(F) \between \mathrm{Multomb}(G)$ implique $F \between G$, --- même si $F$ et $G$ sont \uncertain{réduits} à deux lieux distincts $x$ et $y$ --- auquel cas $\mathrm{Multomb}(X) = \lbrace \lbrace x \rbrace \rbrace$, $\mathrm{Multomb}(Y) = \lbrace \lbrace y \rbrace \rbrace$ sont compatibles et même disjoints, mais $x$ et $y$ ne le sont pas forcément. Mais il sera vrai \ill{} équivalence, dans le cas où $\mathfrak{F}$ est un atelier ensembliste.\note{« $\mathrm{Multomb}(X)$ », « $\mathrm{Multomb}(Y)$ » sur la page, pour les figures $F$, $G$ réduites aux lieux $x$, $y$}

\marginal{vérifier plus tard}

\textbf{Corollaire 2} Si $F$ et $G$ sont disjointes, alors $\mathrm{Multomb}(F)$ et $\mathrm{Multomb}(G)$ aussi.

Cela provient formellement du corollaire 1 et de la formule
\[
  \mathrm{Multomb}(\emptyset_{\mathfrak{F}}) = \emptyset \quad (\text{figure vide dans } \mathcal{M}) \tag{1.40}
\]

\textbf{Proposition} Soient $F, G \in \mathfrak{F}$. Alors $F \ll G \Longleftrightarrow \mathrm{Multomb}(F) \ll \mathrm{Multomb}(G)$.

\page{25}
D'abord une observation. Soit $F \in \mathfrak{F}$, d'où
\[
  \mathrm{Omb}(F) = |\mathrm{Multomb}\, F| \overset{\mathrm{def}}{=} \bigcup_{X \in \widetilde{F}} \mathrm{Omb}(X)
\]
On a une application canonique
\[
  \mathrm{Omb}(F) \xrightarrow{\ \varphi_F\ } \widetilde{F}
\]
\marginal{NB. $\varphi_F$ induit $\mathrm{id}_{\widetilde{F}}$ sur $\widetilde{F} \subset \mathrm{Omb}(F)$}
associant à tout $Y \in \mathrm{Omb}\, F$, i.e. $Y \in \mathcal{M}$ tel que \struck{\ill{}} $Y \ll F$, le plus petit élément de $F^{Y}$ i.e. le plus petit $X$ de $\widetilde{F}$ parmi ceux satisfaisant $Y \ll X$ i.e. \struck{\ill{}} $Y \in \mathrm{Omb}\, X$. Cette application est celle qui définit la relation de préordre de $|\mathrm{Multomb}(F)|$ associée à la figure $\mathrm{Multomb}(F)$.

Ceci dit, dire que $\underbrace{\mathrm{Multomb}(G)}_{\Psi} \ll \underbrace{\mathrm{Multomb}\, F}_{\Phi}$, \add{pour $G, F \in \mathfrak{F}$ données,} signifie qu'on a une inclusion
\[
  |\Psi| \subset |\Phi| \qquad \text{i.e.} \qquad \mathrm{Omb}\, G \subset \mathrm{Omb}\, F,
\]
et une application \emph{croissante} \struck{\ill{}}
\[
  f : \widetilde{G} \to \widetilde{F}
\]
rendant commutatif le diagramme \uncertain{(1.42)}\note{le numéro se lit plutôt « (1.48) » ; ses renvois des pages 26 et 27 à « (1.42) » visent ce diagramme}

\begin{tikzcd}
  \mathrm{Omb}\, G \arrow[r, hook] \arrow[d, "\varphi_G"'] & \mathrm{Omb}\, F \arrow[d, "\varphi_F"] \\
  \widetilde{G} \arrow[r, "f"] & \widetilde{F}
\end{tikzcd}

\page{26}
Montrons que cela équivaut à $G \ll F$. En effet, si $G \ll F$, associant à tout $X \in \widetilde{G}$ le plus petit élément $f(X)$ de $\widetilde{F}$ tel que
\[
  X \ll f(X),
\]
on trouve bien une application
\[
  f : \widetilde{G} \to \widetilde{F}
\]
et elle est évidemment croissante : si $X, X' \in \widetilde{G}$,
\[
  X' \leq X \ \text{ i.e. } \ X' \ll X \ (\ll f(X)) \ \text{ implique que } \ X' \ll f(X), \ \text{ donc } \ f(X') \ll f(X),
\]
qui signifie $f(X') \leq f(X)$\note{mise en page de la déduction condensée : il l'écrit en escalier sur trois lignes}. Et il faut vérifier la commutativité de (1.42), ce qui \struck{signifie \ill{}} \ill{}.

\textbf{Corollaire} Soient $X \in \mathcal{M}$, $F, G \in \mathfrak{F}$ tels que $X \ll G \ll F$. Soit \struck{\ill{}} $X'$ le plus petit élément de $G^{X}$, $X''$ le plus petit élément de $F^{X'}$, alors $X''$ est aussi le plus petit élément de $F^{X}$.

Il semble qu'il faille le prendre en axiome, sous forme \ill{} \ill{}, mais en fait équivalente. On va \struck{\ill{}} introduire une notation, pour $X, Y \in \mathcal{M}$
\[
  X \mathrel{\overset{\circ}{\ll}} Y \tag{1.43}
\]
\note{numéro lu « (1.48) », comme page 25 ; la suite des numéros demande (1.43)}
($X$ est \emph{raffinement intérieur} de $Y$) pour \struck{$X \ll Y$, et $X \ll Y$}
\[
  X \in \mathrm{Omb}(Y)^{\circ} = \mathrm{Omb}(Y) \smallsetminus \bigcup_{Y' < Y} \mathrm{Omb}(Y')
\]
i.e. $X \ll Y$, et $X \ll Y' \leq Y$ implique $Y' = Y$.

\page{27}
On doit donc poser

\textbf{At 7}\note{« At » écrit, comme page 21, par-dessus une autre lettre} Si $X, Y, Z \in \mathcal{M}$ sont tels que
\[
  X \mathrel{\overset{\circ}{\ll}} Y \mathrel{\overset{\circ}{\ll}} Z \quad \text{alors} \quad X \mathrel{\overset{\circ}{\ll}} Z
\]
i.e. la relation $\overset{\circ}{\ll}$ est \emph{transitive}.

Ceci prouve donc $\Longrightarrow$ dans la prop. En sens inverse, s'il existe $f$ comme dans (1.42), soit $X \in \widetilde{G}$, \struck{alors $\varphi_G(X)$} de sorte que $\varphi_G(X) = X$ (tautologique), \uncertain{moyennant} (6), \struck{on a donc} soit $f(X) = Y$, on a donc $Y = \varphi_F(X)$ i.e. $Y$ est le plus petit élément de $\widetilde{F}$ \struck{contenant $X$} tel que $X \ll Y$. Cela \struck{signifie} \struck{Donc} implique donc que pour tout $X \in \widetilde{G}$, existe $Y \in \widetilde{F}$ telle que $X \ll Y$, i.e. (C 4) $\widetilde{G} \ll \widetilde{F}$, cqfd. On n'a pas utilisé le fait que $f$ soit croissante, donc

\textbf{Corollaire \uncertain{2} de la prop.} \struck{Les conditions} \add{Pour que} $\mathrm{Multomb}(F) \ll \mathrm{Multomb}(G)$ i.e. que $F \ll G$, \add{il faut et il suffit} que $\mathrm{Omb}\, F \subset \mathrm{Omb}\, G$, et que \struck{dans} la relation d'équivalence sur $\mathrm{Omb}\, F$ définie par $\varphi_F$ soit plus fine que la relation d'équiv. induite \struck{par} \add{sur} $\mathrm{Omb}\, G$.

\marginal{Ce ne sera \ill{} vrai, il faut \ill{} $R_F$ \ill{} petites}\note{note oblique dans la marge gauche, en face du corollaire, lue en partie}

\page{28}
\textbf{Scholie} Ainsi, via $F \mapsto \mathrm{Multomb}(F)$, $(\mathfrak{F}, \leq, \ll)$ s'interprète comme un « sous-atelier » de $\mathrm{Fig}(\mathcal{M})$.\note{le $\leq$ du triplet est surchargé, comme page 38} Donc la donnée de l'atelier $(\mathfrak{F}, \leq, \ll)$ équivaut à celle d'un ens. $\mathcal{M}$, et d'une partie
\[
  \mathfrak{F}_{\mathcal{M}} \subset \mathrm{Fig}(\mathcal{M}) \subset \mathfrak{P}(\mathfrak{P}(\mathcal{M})) \tag{1.44}
\]
\note{entre $\mathfrak{F}_{\mathcal{M}}$ et le premier $\subset$, des signes lourdement biffés}
(donc $\mathfrak{F}_{\mathcal{M}} \in \mathfrak{P}(\mathfrak{P}(\mathfrak{P}(\mathcal{M})))$) satisfaisant certaines conditions (qui doivent traduire C1 à C7).

Du coup, la notion de multiombre \add{(générale)} reprend un plus grand intérêt, et je me demande s'il ne faudrait pas une notation \struck{plus} \ill{} et un nom plus lapidaire.

Je vais expliciter les conditions sur cet ens. de figures ensemblistes dans $\mathcal{M}$. Je désignerai par At*1 etc. les \ill{} en termes de $(\mathcal{M}, \mathfrak{F}_{\mathcal{M}} \in \mathfrak{P}\mathfrak{P}\mathfrak{P}(\mathcal{M}))$.

\struck{\ill{}} Il est entendu que $\mathfrak{F}_{\mathcal{M}} \subset \mathrm{Fig}(\mathcal{M})$, i.e. \struck{At*1} que les éléments de $\mathfrak{F}_{\mathcal{M}}$ sont des ensembles de parties de $\mathcal{M}$ qui sont des \emph{figures} ensemblistes $F$ dans $\mathcal{M}$. \struck{Ainsi} Considérons

\page{29}
\[
  \mathcal{M}^{*} = \bigcup_{\Phi \in \mathfrak{F}_{\mathcal{M}}} \Phi \subset \mathfrak{P}(\mathcal{M}) \quad = \quad \lbrace A \in \mathfrak{P}(\mathcal{M}) \mid \exists\, \Phi \in \mathfrak{F}_{\mathcal{M}} \text{ avec } A \in \Phi \rbrace
\]
i.e. l'ensemble des strates de toutes les figures $\in \mathfrak{F}_{\mathcal{M}}$.\note{au-dessus de $\mathcal{M}^{*}$, ici et plus bas, un petit signe biffé}

\textbf{At*1} $\mathcal{M}^{*}$ est une figure ensembliste dans $\mathcal{M}$, de support $\mathcal{M}$. De plus, le préordre associé \add{(sur $\mathcal{M}$)} est une relation d'\emph{ordre} i.e.
\[
  X, Y \in \mathcal{M},\ X \neq Y \Longrightarrow \exists\, A \in \mathcal{M}^{*} \text{ avec } X \in A,\ Y \notin A \ \text{ ou } \ Y \in A,\ X \notin A
\]

Ainsi, l'application
\[
  X \longmapsto \mathcal{M}_{\ll X} \overset{\mathrm{def}}{=} \mathrm{Omb}(X), \qquad (*) \quad \mathcal{M} \longrightarrow \mathfrak{P}(\mathcal{M})
\]
induit une \emph{bijection} de $\mathcal{M}$ sur $\mathcal{M}^{*}$, transformant $\ll$ en la relation d'inclusion. Les $A \in \mathcal{M}^{*}$ sont les parties fermées \emph{irréductibles} de $\mathcal{M}$ admettant un pt générique.

\struck{\ill{}} Soit $\Phi \in \mathfrak{F}_{\mathcal{M}}$ (donc $\Phi \subset \mathfrak{P}(\mathcal{M})$ et par définition même $\Phi \subset \mathcal{M}^{*} \subset \mathfrak{P}(\mathcal{M})$), donc $\Phi$ est

\page{30}
l'image par \struck{(1.39)} (*)
\[
  \mathcal{M} \xrightarrow{\ \sim\ } \mathcal{M}^{*} \subset \mathfrak{P}(\mathcal{M})
\]
d'une unique \add{partie} $F \subset \mathcal{M}$, à savoir que $\Phi = \widetilde{F}$, i.e.
\[
  \Phi = \lbrace \mathrm{Omb}(X) \mid X \in F \rbrace, \qquad \mathrm{Omb}(X) = \mathcal{M}_{\ll X}
\]
Ainsi, on peut aussi interpréter, grâce à $\ll$, $\mathfrak{F}_{\mathcal{M}}$ en termes d'une \struck{partie} \add{ens. de parties} $\widetilde{\mathfrak{F}}$ de $\mathcal{M}$

\begin{tikzcd}
  \mathfrak{F}_{\mathcal{M}} \arrow[r, "\sim"] \arrow[d, no head, "\cap" description] & \widetilde{\mathfrak{F}} \arrow[d, no head, "\cap" description] \\
  \mathfrak{P}(\mathfrak{P}(\mathcal{M})) & \mathfrak{P}(\mathcal{M})
\end{tikzcd}

\note{les inclusions verticales sont notées par un $\cap$ sous chaque terme}

Ceci dit, \struck{on demande \ill{} que} \add{ceux \ill{}} $\widetilde{\mathfrak{F}} \subset \mathfrak{P}(\mathcal{M})$, \struck{satisfaisant \ill{}} \ill{} satisfont C1, C2, C3, \struck{la relation d'inclusion \ill{}} \ill{} par $\mathcal{M}$ s'identifie à l'un des \ill{} objets \add{\ill{}} \struck{\ill{}} l'ens. des parties $\widetilde{\mathfrak{F}} \subset \mathfrak{P}(\mathcal{M})$ \add{(\ill{} d'abord que)} \struck{soit} \emph{préadmissible}, i.e. pour tout $X \in \mathcal{M}$, il existe un plus petit élément de $\widetilde{\mathfrak{F}}$ qui le contienne. En d'autres termes,\note{paragraphe très corrigé : plusieurs lignes biffées et des ajouts interlinéaires, lus en partie ; « Axiome » en tête de ligne est biffé}

\marginal{cf GF IV page 54 (cor. 2) et p. 55 (variante), id pour C3, p. 58}\note{renvoi, dans la marge gauche, à la première mouture (chapitre IV, dossier 156-4), selon sa pagination ; c'est là que les conditions C1 à C3 étaient établies}

\textbf{At*2}\note{le sigle porte en indice deux petits ronds, peut-être un « $\circ$ » doublé} Pour toute \struck{$\Phi \in \mathfrak{F}_{\mathcal{M}}$} $\mathcal{M}$-figure ensembliste $\Phi \in \mathfrak{F}_{\mathcal{M}}$, et pour toute \struck{strate de} $A$ \add{$[= \mathrm{Omb}(X)\,!]$} de $\Phi$ (i.e. $A \in \Phi$), il existe une \emph{plus petite figure} $\Phi_A$ \struck{\ill{}} $\in \mathfrak{F}_{\mathcal{M}}$ contenant $A$ (plus petite au sens de l'\emph{inclusion} entre parties de $\mathfrak{P}(\mathcal{M})$, ou de $\mathcal{M}^{*} \subset \mathfrak{P}(\mathcal{M})$ --- figures i.e. entre parties qui contiennent $A$ comme élément (laquelle $\Phi_A$ sera donc $\subset \Phi$)).

\marginal{c'est la multistrate de figure $A$}

\textbf{At*3} Si $\Phi \in \mathfrak{F}_{\mathcal{M}}$, et si $\Phi_i \in \mathfrak{F}_{\mathcal{M}}$, $\Phi_i \subset \Phi$, alors $\bigcup \Phi_i \in \mathfrak{F}_{\mathcal{M}}$ (stabilité par réunions majorées dans $\mathfrak{F}_{\mathcal{M}}$)

\marginal{C'est pour $\mathfrak{F}_{\mathcal{M}}$ C 1}

\page{31}
et de plus on a $\emptyset \in \mathfrak{F}_{\mathcal{M}}$ (en outre $\mathfrak{F}_{\mathcal{M}} \neq \emptyset$).

Enfin, on voudra aussi

\textbf{At*4} Si $\Phi, \Psi \in \mathfrak{F}_{\mathcal{M}}$, pour que $\Phi \cup \Psi$ soit \struck{élément} dans $\mathfrak{F}_{\mathcal{M}}$ \add{(i.e., par At*3, $\Phi$, $\Psi$ majorés dans $\mathfrak{F}_{\mathcal{M}}$ \ill{})}, il faut et il suffit que pour tout $A \in \Phi$, $B \in \Psi$, $\Phi_A \cup \Psi_B$ soit dans $\mathfrak{F}_{\mathcal{M}}$.

Je voudrais revenir sur ce $\Phi_A$. En fait, $\Phi_A$ consiste en l'ensemble des sous-strates de $A$ dans $\Phi$, donc At*2 peut se formuler sous la forme plus compréhensible

\textbf{At*2.} Soient $\Phi$ et $\Psi$ deux $\mathcal{M}$-figures $\in \mathfrak{F}_{\mathcal{M}}$, et soit $A \in \Phi \cap \Psi$ une strate commune. Alors la multistrate associée : $A$ dans $\Phi$, et dans $\Psi$, est la même, et elle est $\in \mathfrak{F}_{\mathcal{M}}$.

\textbf{Corollaire} Si $\Phi, \Psi \in \mathfrak{F}_{\mathcal{M}}$ avec $\Phi \subset \Psi$, alors $\Phi$ est fermée dans $\Psi$, donc est une sous-\emph{figure} ($\Phi \leq \Psi$).

\page{32}
[Il s'ensuit qu'on a dans $\mathcal{M} \simeq \mathcal{M}^{*}$ une relation d'ordre $\leq$, avec
\[
  X \leq Y \Longleftrightarrow \Phi_{\mathrm{Omb}(X)} \subset \Phi_{\mathrm{Omb}(Y)}, \quad \text{i.e.} \quad \mathrm{Omb}(X) \in \Phi_{\mathrm{Omb}(Y)} \ \overset{\mathrm{def}}{=}\ \mathrm{Multomb}(Y)
\]
\note{avant $\Phi_{\mathrm{Omb}(X)}$, « $\mathrm{Omb}(X)$ » est biffé}
de telle façon que les \struck{éléments de} $F \in \widetilde{\mathfrak{F}}$ soient des parties fermées de $\mathcal{M}$. Associant à tout $X \in \mathcal{M}$ la partie
\[
  \widetilde{X} = \lbrace Y \in \mathcal{M} \mid Y \leq X \rbrace = \mathcal{M}_{\leq X}
\]
\note{dans l'accolade, après « $Y \leq X$ » : « i.e. $\mathrm{Omb}(Y) \in \Phi_{\mathrm{Omb}(X)}$ »}
on trouve une application injective
\[
  \mathcal{M} \hookrightarrow \widetilde{\mathfrak{F}} \quad (\simeq \mathfrak{F}_{\mathcal{M}})
\]
qui induit un iso entre $\mathcal{M}$ et l'ens. des éléments str. irréductibles de $\widetilde{\mathfrak{F}}$ (ordonné par inclusion, relation $\leq$).]

\struck{D}

\textbf{Théorème-Scholie} La catégorie des \add{(\uncertain{quasi-})} ateliers $(\mathfrak{F}, \leq, \ll)$ satisfaisant At 1 \ldots\ At 7, équivaut à celle des $(\mathcal{M}, \mathfrak{F}_{\mathcal{M}})$ (avec $\mathfrak{F}_{\mathcal{M}} \subset \mathrm{Fig}(\mathcal{M})$) satisfaisant les conditions suivantes, \ill{} :

\textbf{At*1} $\mathfrak{F}_{\mathcal{M}}$ \add{contient $\emptyset$ (figure vide), et} \struck{non vide, et} (stable par \struck{réunions majorées} \add{réunions majorées dans $\mathfrak{F}_{\mathcal{M}}$})

\struck{Cette} [La stabilité par \ill{}, \ill{} la condition suivante, \ill{} celle de \struck{\ill{}} stabilité par \ill{} : une sous-\struck{figure} \ill{} \emph{figure})\note{fin de page très corrigée, lue en partie ; la parenthèse n'est pas refermée sur la page}

\page{33}
\textbf{At*2} Si $\Phi, \Psi \in \mathfrak{F}_{\mathcal{M}}$, et $A \in \Phi \cap \Psi$, alors $\Phi_A = \Psi_A$, i.e. la multistrate définie par $A$, soit \struck{\ill{}} dans $\Phi$, soit dans $\Psi$, est la même.

On l'appellera \ill{} \struck{$A$}. « Toute strate \uncertain{porte} \uncertain{canoniquement} une structure de multistrate », indépendamment des $\Phi \in \mathfrak{F}_{\mathcal{M}}$ dont c'est une strate.

\marginal{At*1 et At*2 \ill{} (C1 et C2 \ill{}) $(\mathfrak{F}_{\mathcal{M}}, \leq)$}\note{note oblique dans la marge gauche, en haut de page, lue en partie}

\textbf{At*4}\note{le chiffre ressemble à un « 9 » ; la liste de la page 31 demande At*4} Si $\Phi, \Psi \in \mathfrak{F}_{\mathcal{M}}$ sont telles que $\forall\, A \in \Phi$, $B \in \Psi$, \struck{\ill{}} on ait $\Phi_A \cup \Psi_B \in \mathfrak{F}_{\mathcal{M}}$, alors $\Phi \cup \Psi \in \mathfrak{F}_{\mathcal{M}}$.

\marginal{NB il suffit que $\Phi_A$, $\Psi_B$ soient compatibles comme figures ensemblistes !}

\textbf{At*5} L'ensemble \struck{\ill{}} $\mathcal{M}^{*}$ des \add{$A$} strates \struck{\ill{}} \add{de toutes} les $\Phi \in \mathfrak{F}_{\mathcal{M}}$ est une figure dans $\mathcal{M}$ de support $\mathcal{M}$. De plus, \add{pour} $X, Y \in \mathcal{M}$, $X \neq Y$, existe $A \in \mathcal{M}^{*}$ telle que $Y \in A$, $X \notin A$ ou $X \in A$, $Y \notin A$ ($\mathcal{M}^{*}$ « sépare les points » de $\mathcal{M}$).

\textbf{At*3} Soient $A, B$ deux strates \add{\ill{}} (i.e. des strates de $\Phi, \Psi \in \mathfrak{F}_{\mathcal{M}}$), telles que $A \subset B$, alors $\widetilde{A} \ll \widetilde{B}$.\note{At*3 est entouré d'une accolade, et une flèche tracée dans la marge le fait remonter avant At*4, sous At*2}

\marginal{i.e. tout $X \in \mathcal{M}$ \ill{} strate \ill{} \ill{} ; NB On \ill{} que At*2 et \ill{} At*3 \ill{} $\Phi \in \mathfrak{F}_{\mathcal{M}}$, $B \in \Psi$, $A \subset B$ \ill{} $\Phi_A \ldots \Psi_B$ \ill{}}\note{longues notes obliques dans la marge gauche de la moitié inférieure, lues par fragments}

J'espère que je n'ai pas oublié d'axiome !

\page{34}
\note{en tête de page, souligné : « 20 juin » — date de sa main}

\emph{20 juin} Comme
\[
  \mathrm{Fig}(\mathcal{M}) \longrightarrow \mathrm{Fig}(\mathcal{L})
\]
induit par l'application
\[
  \mathcal{L} \longrightarrow \mathcal{M}
\]
est croissant pour $\leq$ et $\ll$, il s'ensuit que l'application
\[
  \mathfrak{F} \longrightarrow \mathrm{Fig}(\mathcal{L}), \qquad F \longmapsto \mathrm{multomb}(F)
\]
(\emph{petites} multiombres) est également croissante pour $\leq$ et $\ll$, et commute aux Sup pour $\leq$, et à Inf pour $\leq$. Mais elle n'est plus nécessairement injective. [Elle transforme sous-figures en sous-figures]

\textbf{Exemple 1} A) Soit $E$ espace topologique, \struck{$\ill{}$} $\mathfrak{F}$ l'ensemble des \struck{figures ensemblistes} \struck{paires} triples
\[
  F = (X, \Sigma, \Phi)
\]
où
\begin{itemize}
  \item $X$ est une partie fermée de $E$,
  \item $\Sigma$ est une structure lin. par morceaux (\ill{}, \ill{} donnée) sur $X$ \uncertain{modérée} \ill{},
  \item $\Phi$ est une figure ens. sur $X$, de support $X$, dont les strates sont des sous-ens. fermés lin. par morceaux.
\end{itemize}
\[
  (X', \Sigma', \Phi') \leq (X, \Sigma, \Phi) \overset{\mathrm{df}}{\Longleftrightarrow} \begin{cases} \text{a) } X' \subset X,\ \text{sous-espace lin. par morceaux pour } \Sigma \\ \text{b) } \Sigma' \text{ induite par } \Sigma \\ \text{c) } \Phi' \leq \Phi \end{cases}
\]
\[
  (X', \Sigma', \Phi') \ll (X, \Sigma, \Phi) \overset{\mathrm{df}}{\Longleftrightarrow} \begin{cases} \text{a) et b) comme dessus} \\ \text{c) } \Phi' \ll \Phi \end{cases}
\]
On \ill{} \ill{} l'ens. $\mathcal{L}$ s'identifie à $E$, et l'application $\mathfrak{F} \to \mathrm{Fig}(E) = \mathrm{Fig}(E)$ est l'application $F = (X, \Sigma, \Phi) \mapsto \Phi$. Elle n'est visiblement pas injective (en général p. ex. si $E$ contient \ill{} homéo à $[0,1]$).

B) \emph{Variante} $\mathcal{M} =$ ens. des \struck{parties de} \add{couples} $X = (A, \sigma)$ avec
\begin{itemize}
  \item $A$ partie \struck{fermée} de $E$
  \item $\sigma$ une structure de \emph{simplexe affine} sur $A$, d'où figure ensembliste $\mathfrak{L}_{\sigma}$ avec $|\mathfrak{L}_{\sigma}| = A$.
\end{itemize}
On suppose l'inclusion $A \subset E$ continue.

\page{35}
On pose
\[
  \begin{aligned}
    (A', \sigma') \leq (A, \sigma) &\overset{\mathrm{df}}{\Longleftrightarrow} (A', \sigma') \text{ sous-simplexe (face) de } (A, \sigma) \\
    (A', \sigma') \ll (A, \sigma) &\overset{\mathrm{df}}{\Longleftrightarrow} \mathfrak{L}_{\sigma'} \ll \mathfrak{L}_{\sigma} \\
    (A, \sigma) \between (A', \sigma') &\Longleftrightarrow A \cap A' \text{ est vide ou une face commune de } (X, \sigma) \text{ et } (X', \sigma')
  \end{aligned}
\]
\note{« de $(A,\sigma)$ », « vide ou » sont ajoutés au-dessus de la ligne ; la dernière ligne écrit $(X,\sigma)$, $(X',\sigma')$ pour $(A,\sigma)$, $(A',\sigma')$}

$\mathfrak{F} =$ ens. des parties finies (p. ex.) de $\mathcal{M}$, fermées pour $\leq$, et formées de simplexes $(A, \sigma)$ deux à deux compatibles.

(NB au lieu de « finies », on pourrait prendre aussi « localement finies »)

Ici encore, $\mathcal{L} \simeq E$, et ici
\[
  \mathcal{M} \longrightarrow \mathrm{Mulstrat}(\mathcal{L}) \subset \mathrm{Fig}(\mathcal{L})
\]
n'est en général pas injectif (par p. ex. si $E$ contient une partie homéomorphe à $[0,1]$ !)

\textbf{Exemple 2} Considérons un « pseudodisque » $(B, \delta B, \ldots)$ dans esp. top. $E$. \struck{$B \supset B$. Alors un} Alors, \struck{la donnée (\ill{})} pour une définition convenable de l'ens. des $\mathfrak{F}$ de (pseudo)-figures de $E$, indiquant le type \struck{topologique} \add{combinatoire} (disque avec pt marqué) une figure (\ill{}) de ce type, correspondant à $(B, \delta B)$, \ill{} par la seule \ill{} de la partie (contractile) \ill{} de $B, \delta B$ \ill{} « \ill{} » du disque, mais il faut de plus se donner une « structure \ill{} » \ill{} au bord \ldots\note{avant « de ce type », un petit cercle contenant un signe, sans doute le dessin du disque à point marqué ; l'exemple est écrit vite et lu en partie}

Ceci étant, il est \add{\uncertain{sans doute}} illusoire d'espérer que
\[
  \begin{aligned}
    \mathrm{multomb}(F) \leq \mathrm{multomb}(F') &\Longrightarrow F \leq F' \\
    \mathrm{multomb}(F) \ll \mathrm{multomb}(F') &\Longrightarrow F \ll F'
  \end{aligned}
\]
\note{au début de la première ligne, un premier départ (« \ill{} $\mathrm{Mult}\,\mathrm{multomb}(F)$ ») est biffé}

\page{36}
et en fait, le contraire est vrai dans les exemples 1. A) et B).

\textbf{Définition} On dit que $\mathfrak{F}$ est un \emph{atelier ensembliste} si pour $F, G \in \mathfrak{F}$\add{, $\mathfrak{F} \to \mathrm{Fig}(\mathcal{L})$ injectif}
\begin{itemize}
  \item[a)] $\mathrm{multomb}\, F \leq \mathrm{multomb}\, G \Longrightarrow F \leq G$
  \item[b)] $\mathrm{multomb}\, F \ll \mathrm{multomb}\, G \Longrightarrow F \ll G$
  \item[c)] $\mathrm{multomb}\, F \between \mathrm{multomb}\, G \Longrightarrow F \between G$
\end{itemize}
\struck{$\mathrm{multomb}(F)$ élémentaire $\Longleftrightarrow$} (\ill{} pas \ill{} question \ill{} : \ill{} conditions), de telle sorte que
\[
  \mathfrak{F} \longrightarrow \mathrm{Fig}(\mathcal{L})
\]
identifie $\mathfrak{F}$, avec les structures $\leq$, $\ll$, $\between$, à un « sous-atelier » de $\mathrm{Fig}(\mathcal{L})$.

\marginal{\ill{} élémentaire \ill{} $\mathrm{multomb}(F)$ élémentaire $\Longleftrightarrow$ $F$ élémentaire}\note{note oblique dans la marge gauche, en face de c), lue en partie}

\struck{\ill{}} Je voudrais préciser ce point. \add{D'abord}

\textbf{Définition} Soient $(\mathfrak{F}, \leq, \ll)$ et $(\mathfrak{F}', \leq, \ll)$ deux ateliers, $f = f_{\mathrm{fig}} : \mathfrak{F} \to \mathfrak{F}'$ une application. On dit que $f$ est un \emph{morphisme d'ateliers} si
\begin{itemize}
  \item[a)] $f$ croissante pour $\leq$ et $\ll$ (donc $\emptyset_{\mathfrak{F}} \mapsto \emptyset_{\mathfrak{F}'}$)
  \item[b)] \struck{$f$ \ill{}} $f$ commute aux \emph{Sup} \ill{} pour $\leq$ \add{et aux \ill{} majorés} \struck{\ill{} figures \ill{}}
\end{itemize}
\struck{$f : \mathcal{M} \to \mathcal{M}'$} On dit que $f$ est « compatible aux strates » si \struck{telle que $f$ croissante pour} $f$ satisfait
\begin{itemize}
  \item[c)] $f$ transforme multistrates (i.e. fig. él.) en multistrates (fig. él.)
\end{itemize}
Alors la donnée d'une telle $f$ équivaut à celle de
\[
  f_{\mathrm{mult}} : \mathcal{M} \longrightarrow \mathcal{M}'
\]

\marginal{NB \ill{} $\mathcal{F} \subset \mathfrak{F}$ \ill{} $\ll$ \ill{} \ill{} $F \leq F'$ \ill{} \ill{} Dans le \ill{} des ensemblistes, \ill{} croissante \ill{} également \ldots}\note{longues notes obliques dans la marge gauche de la moitié inférieure, dont une ouverte par « b) Dans le » ; lues par fragments}

\page{37}
satisfaisant
\begin{itemize}
  \item[a)] $f_{\mathrm{mult}}$ croissante pour $\leq$ et $\ll$
  \item[b)] \struck{$f_{\mathrm{m}}$} pour tout $F \subset \mathcal{M}$ provenant d'une figure, $f_{\mathrm{m}}(F)$ est une figure
\end{itemize}

\textbf{NB} Pour que $f$ soit \ill{}, il faut et il suffit que $f_{\mathrm{m}}$ le soit. Pour que ce soit un iso sur un \ill{} \ill{}, il faut et il suffit que $f$ soit un plongement pour $\leq$ et $\ll$, et que l'image soit \ill{}-fermée.\note{NB encadré, écrit serré entre les lignes, lu en partie}

On appelle \emph{sous-atelier} d'un atelier $\mathfrak{F}$, une \struck{de} partie $\mathfrak{F}'$ de $\mathfrak{F}$, telle que, munie des relations d'ordre $\leq_{\mathfrak{F}'}$, $\ll_{\mathfrak{F}}$ induites par celles de $\mathfrak{F}$, ce soit un atelier, et que $\mathfrak{F}' \hookrightarrow \mathfrak{F}$ soit un morphisme d'ateliers, i.e. que
\begin{itemize}
  \item[1°)] si $F_i \in \mathfrak{F}'$ sont majorés dans $\mathfrak{F}'$ \add{par $F$}, alors $\operatorname{Sup}_{i \in I} F_i \in \mathfrak{F}'$ (stabilité par Sup majorés)
  \item[2°)] sous les \struck{\ill{}} hypothèses, $\bigcap F_i \in \mathfrak{F}'$
\end{itemize}
\ill{} je me demande s'il ne serait pas plus raisonnable, \ill{} de 1°, et 2°, de poser \struck{\ill{}}
\begin{itemize}
  \item[(*)] si $F' \in \mathfrak{F}'$, \struck{et $F \in \mathfrak{F}$ est tel que \ill{}} \ill{} $F'$ \ill{} dans $\mathfrak{F}'$.
\end{itemize}
\note{« (*) » est entouré sur la page}
Il faut vérifier \add{\ill{}} que cette condition suffit pour que $(\mathfrak{F}', \leq_{\mathfrak{F}'}, \ll_{\mathfrak{F}'})$ soit bien un atelier \add{(\emph{strict})}\ldots\ Les sous-ateliers correspondant aux sous-ens. $\mathcal{M}'$ de $\mathcal{M}$ \uncertain{fermés} pour $\leq$ \ill{}, \ill{} $\mathcal{M}'$ s'identifie à l'un des \ill{}. Donc, si $\mathfrak{F}$ est dit ensembliste, i.e. l'application
\[
  \mathrm{multomb} : \mathfrak{F} \longrightarrow \mathrm{Fig}(\mathcal{L})
\]
le « réalise » comme un sous-\ill{} de $\mathrm{Fig}\, \mathcal{L}$ --- en fait, on aimerait que ce soit un sous-atelier \emph{strict}\,$^{*}$, et on voudra \struck{que} en particulier que $\mathcal{M}$ se réalise comme un ens. de figures (multistrates ensemblistes) élémentaires \emph{dans} $\mathcal{L}$. Ainsi, dans la définition de la page précédente, il faudrait \ill{} poser

\marginal{Mais les sous-ateliers \ill{} pas stricts \ill{} $f(\mathcal{M}) \subset \mathcal{M}'$ \ill{} ; Mais on \ill{} $F, F' \in \mathfrak{F}'$ \ill{} $F \vee F'$ \ill{} ; $^{*}$ c'est-à-dire \ill{} \ill{} stabilité \ill{} $F \vee$ \ill{}}\note{notes obliques dans la marge gauche, sur toute la hauteur de la page, lues par fragments ; l'astérisque de « \emph{strict} » y renvoie}

\page{38}
\note{toute la page est barrée d'un long trait oblique, du haut au bas du feuillet ; elle est donnée en entier}

des conditions : $\mathfrak{F}$ stable dans $\mathrm{Fig}(\mathcal{L})$ par passage à une sous-figure, et par $\operatorname{Sup}_{\leq}$ (en demandant \ill{} \ill{}).

Je reprends ces définitions

\textbf{Définition} Soit $(\mathfrak{F}, \leq, \ll)$ un atelier At 1 -- At 7 (on verra plus loin des raisons ultérieures pour \ill{}). Une \emph{sous-atelier} \struck{(\ill{})} est une partie \add{non vide} $\mathfrak{F}'$ de $\mathfrak{F}$, fermée pour $\leq$ (i.e. telle que toute sous-figure d'une figure dans $\mathfrak{F}'$ est dans $\mathfrak{F}'$) \add{et stable par Sup majorés dans $\mathfrak{F}'$}. On dit que $\mathfrak{F}'$ est \emph{spécieux} \add{de plus} s'il est stable par Sup quelconques (\ill{} supposés majorés dans $\mathfrak{F}'$).

\marginal{non vide i.e. $\emptyset_{\mathfrak{F}} \in \mathfrak{F}'$}

Ces notions \struck{\ill{}} ne dépendent que de $\leq$. Mais on munira bien sûr $\mathfrak{F}'$ des deux relations $\leq$, $\ll$ induites.

Soit
\[
  \mathcal{M}' = \mathfrak{F}' \cap \mathcal{M}
\]
alors on voit que les objets de $\mathcal{M}'$ sont les objets (strictement) irréductibles de $\mathfrak{F}'$ pour $\leq$ (cf GF IV, p. 53). En effet, si \struck{$X \in \mathfrak{F}$,} \struck{\ill{} $X = \operatorname{Sup} F_i$ dans $\mathfrak{F}$,} \struck{on a aussi $X = \operatorname{Sup} F_i$ dans $\mathfrak{F}$ (\ill{})} \struck{par hyp. donc $\exists\, i$ avec $X = F_i$, donc $X$ irr. dans $\mathfrak{F}'$} \struck{Inversement, si $X \in \mathfrak{F}'$ est irr. dans $\mathfrak{F}$,} les familles $\lbrace F_i \rbrace$ dans $\mathfrak{F}'$, \add{majorées par $X$,} sont les mêmes \struck{\ill{}}. Comme $\mathfrak{F}'$ \struck{est stable par Sup}, dire que pour une telle famille, $X = \operatorname{Sup} F_i$ dans $\mathfrak{F}'$ revient au même. Donc dire que pour \struck{\ill{}} toute telle famille, $\exists\, i$ avec $X = F_i$, est idem pour $\mathfrak{F}$ ou $\mathfrak{F}'$.

\struck{Il faudrait vérifier.}

Dans le cas d'un atelier spécieux on \ill{}

\page{39}
\textbf{Définition} Soit $(\mathfrak{F}, \leq)$ un ens. ordonné \add{avec Sup majorés et plus petit élément $\emptyset_{\mathfrak{F}}$}. Une partie \struck{\ill{}} fermée $\mathfrak{F}'$ de $\mathfrak{F}$ est dite \emph{ample} si $\mathfrak{F}'$ est stable \add{et $\mathfrak{F}' \neq \emptyset$ ($\Leftrightarrow \emptyset_{\mathfrak{F}} \in \mathfrak{F}'$)} par Sup. majorés \emph{dans} $\mathfrak{F}'$. On dit qu'elle est \emph{spécieuse} si elle est \struck{stable par} $=$ stable par Sup majorés dans $\mathfrak{F}$.\note{cette définition est barrée de plusieurs traits obliques}

Soit $(\mathfrak{F}, \leq)$ un ens. ordonné \struck{stable par} avec Sup. majorés et plus petit élément $\emptyset_{\mathfrak{F}}$ (i.e. satisfaisant At 1). Soit $\mathfrak{F}'$ une partie \emph{fermée}. Alors \struck{dans} $\mathfrak{F}'$ est non vide ssi $\emptyset_{\mathfrak{F}} \in \mathfrak{F}'$. D'autre part, si $(F_i)$ est une famille dans $\mathfrak{F}'$ qui est majorée dans $\mathfrak{F}'$ par $G$, alors son sup dans $\mathfrak{F}$, soit $F$, est $\leq G$, donc $\in \mathfrak{F}'$, et c'est son sup dans $\mathfrak{F}'$. Donc $\mathfrak{F}'$ \add{(il satisfait donc At 1)} est stable par Sup. majorés dans $\mathfrak{F}'$, et $\mathfrak{F}' \hookrightarrow \mathfrak{F}$ commute aux Sup.

Soit $\mathcal{M}$ l'ens. des éléments (str.) irréductibles de $\mathfrak{F}$, et
\[
  \mathcal{M}' = \mathcal{M} \cap \mathfrak{F}' .
\]
Alors, $\mathcal{M}'$ est l'ens. des objets (str.) irréd. de $\mathfrak{F}'$. En d'autres termes, si $F \in \mathfrak{F}'$, pour que $F$ soit irréd. dans $\mathfrak{F}'$, il faut et il suffit qu'il le soit dans $\mathfrak{F}$. C'est immédiat, en termes des définitions : les familles $(F_i)$ d'objets majorés par $F$, que ce soit dans $\mathfrak{F}$ ou $\mathfrak{F}'$, sont les mêmes, et leurs Sup dans $\mathfrak{F}$ ou $\mathfrak{F}'$ itou, et donc celles pour lesquelles $F = \operatorname{Sup} F_i$ sont les mêmes, et dire que pour

\page{40}
telle famille, $\exists\, i$ t.q. $F_i = F$, est \add{donc} pareil dans $\mathfrak{F}$ et $\mathfrak{F}'$.)

Supposons que $\mathfrak{F}$ satisfasse At 2 :
\[
  \forall F \in \mathfrak{F}, \quad F = \operatorname*{Sup}_{X \in \widetilde{F}} X \quad \text{où} \quad \widetilde{F} = \lbrace X \in \mathcal{M} \mid X \leq F \rbrace = \lbrace X \in \mathfrak{F},\ X \trianglelefteq F \rbrace
\]
Alors $\mathfrak{F}'$ y satisfait aussi, car si $F \in \mathfrak{F}'$, $\widetilde{F}$ au sens de $\mathfrak{F}$ ou de $\mathfrak{F}'$ c'est pareil, et le sup dans $\mathfrak{F}$ est un sup dans $\mathfrak{F}'$.

\textbf{Définition} La partie fermée $\mathfrak{F}'$ \struck{\ill{}} de $\mathfrak{F}$ \add{\ill{}} est dite \emph{spécieuse} si elle est stable par Sup quelconques (de familles \emph{majorées dans} $\mathfrak{F}$).

\marginal{On dira \ill{} \emph{ample} \ill{} \ill{} stable par \ill{} ; Remarquons \ill{} \ill{} « spécieuse » \ill{}}\note{notes obliques dans la marge gauche, lues par fragments}

\textbf{Exemple} Soit $\mathfrak{F}$ l'ens. des parties d'un ens. $I$ \add{$= \mathfrak{P}(I)$}, avec la relation d'inclusion. Satisfait At 1, At 2, At 3, $\mathcal{M} = \lbrace \lbrace i \rbrace \mid i \in I \rbrace$ (comme \struck{\ill{}} \ill{} discret), \ill{} donner une $\mathfrak{F}'$ fermée $=$ donner une structure de \add{\ill{}} \emph{complexe} simplicial sur $I$, du moins si $I$ est fini (sinon, c'est vrai à condition de se restreindre aux $\mathfrak{F}'$ formés d'ens. finis \ldots). Dire que $\mathfrak{F}'$ est spécieux, signifie que \struck{\ill{}} la réunion de simplexes est un simplexe --- i.e. \ill{} donner une partie $I'$ de $I$, \struck{\ill{}} et $\mathfrak{F}' = \mathfrak{P}(I')$ --- c'est pas du \struck{tout} pareil !

\end{document}
